\chapter{Planning Service: Backend API Integration}
\label{chap:planning-service}

The previous chapters covered client-side functionality: UI state, data management, file I/O, and layout algorithms. This chapter introduces the \textbf{Planning Service}, which extends the application's capabilities by communicating with a powerful backend server for simulation and analysis tasks.

\section{Problem Statement}
\label{sec:planning-problem}

Some operations are too complex or resource-intensive for browser-based execution:

\begin{itemize}
    \item \textbf{Long simulations}: Running thousands of firing steps
    \item \textbf{Stochastic analysis}: Monte Carlo simulations with multiple runs
    \item \textbf{AI planning}: Computing optimal transition sequences
    \item \textbf{Domain-specific algorithms}: Offshore scheduling optimization
\end{itemize}

These tasks require:
\begin{itemize}
    \item More computing power than browsers provide
    \item Specialized libraries (e.g., Python's \texttt{gspn4py} for stochastic Petri nets)
    \item Server-side state management for long-running operations
\end{itemize}

The \texttt{PlanningService} acts as a \textbf{bridge} between the frontend and backend, sending requests and receiving results via HTTP.

\section{System Architecture}
\label{sec:planning-architecture}

\begin{figure}[htb]
    \centering
    \begin{tikzpicture}[
        node distance=2cm,
        box/.style={rectangle, draw, minimum width=3cm, minimum height=1.2cm, align=center},
        arrow/.style={->, thick},
    ]
        % Frontend
        \node[box, fill=green!10] (fe) {Angular Frontend\\(Browser)};
        
        % Backend
        \node[box, fill=blue!10, right=4cm of fe] (be) {Flask Backend\\(Python Server)};
        
        % Arrows
        \draw[arrow] ([yshift=5pt]fe.east) -- node[above] {\small HTTP POST (PNML)} ([yshift=5pt]be.west);
        \draw[arrow] ([yshift=-5pt]be.west) -- node[below] {\small JSON (Results)} ([yshift=-5pt]fe.east);
        
        % Labels
        \node[below=0.5cm of fe] {\small PlanningService};
        \node[below=0.5cm of be] {\small /simulation-petri-nets/};
        
    \end{tikzpicture}
    \caption{Frontend-backend communication via PlanningService}
    \label{fig:planning-architecture}
\end{figure}

\section{Environment Configuration}
\label{sec:planning-environment}

The backend URL is configured via Angular's environment files:

\begin{lstlisting}[language=TypeScript, caption={Development environment (src/environments/environment.ts)}]
export const environment = {
    production: false,
    backendApiUrl: 'http://localhost:9040/l3s-offshore-2'
};
\end{lstlisting}

\begin{lstlisting}[language=TypeScript, caption={Production environment (src/environments/environment.prod.ts)}]
export const environment = {
    production: true,
    backendApiUrl: 'https://l3s-offshore-plan-2.l3s.uni-hannover.de/l3s-offshore-2'
};
\end{lstlisting}

Angular automatically swaps these files during the build process, ensuring development builds connect to \texttt{localhost} while production builds use the live server.

\section{PlanningService Implementation}
\label{sec:planning-implementation}

\begin{lstlisting}[language=TypeScript, caption={Planning service (src/app/tr-services/planning.service.ts)}]
import { Injectable } from '@angular/core';
import { HttpClient, HttpErrorResponse } from '@angular/common/http';
import { Observable, throwError } from 'rxjs';
import { catchError } from 'rxjs/operators';
import { environment } from '../../environments/environment';

@Injectable({
    providedIn: 'root'
})
export class PlanningService {
    /** Base URL for planning-related endpoints */
    private apiUrl = `${environment.backendApiUrl}/dto/planning`;
    
    /** URL for Petri net simulation endpoints */
    private simpleSimApiUrl = `${environment.backendApiUrl}/simulation-petri-nets/simple-sim`;

    constructor(private http: HttpClient) {}

    // =========================================================================
    // SIMULATION METHODS
    // =========================================================================

    /**
     * Run a Petri net simulation on the backend.
     * Sends the PNML file and receives firing sequence results.
     * 
     * @param pnmlFile - The PNML file to simulate
     * @param runs - Number of simulation runs (for stochastic analysis)
     * @returns Observable with simulation results
     */
    runSimpleSimulation(pnmlFile: File, runs: number = 1): Observable<SimulationResult> {
        // FormData is required for file uploads
        const formData = new FormData();
        formData.append('pnml_model', pnmlFile, pnmlFile.name);
        formData.append('num_runs', runs.toString());

        return this.http.post<SimulationResult>(this.simpleSimApiUrl, formData)
            .pipe(catchError(this.handleError));
    }

    /**
     * Convenience method: run simulation from PNML string.
     * Converts the string to a File object before sending.
     */
    runSimpleSimulationFromString(
        pnmlContent: string, 
        runs: number = 1, 
        fileName: string = 'model.pnml'
    ): Observable<SimulationResult> {
        // Create a File from the string content
        const pnmlFile = new File([pnmlContent], fileName, { type: 'application/xml' });
        return this.runSimpleSimulation(pnmlFile, runs);
    }

    // =========================================================================
    // PARAMETER & CONFIGURATION METHODS
    // =========================================================================

    /**
     * Fetch default parameter values from backend.
     * Used to populate the parameter table with sensible defaults.
     */
    getDefaults(): Observable<ParameterDefaults> {
        return this.http.get<ParameterDefaults>(`${this.apiUrl}/defaults`)
            .pipe(catchError(this.handleError));
    }

    /**
     * Submit planning parameters to backend.
     * Triggers the offshore scheduling optimization.
     */
    runPlanning(planningData: PlanningRequest): Observable<PlanningResult> {
        return this.http.post<PlanningResult>(this.apiUrl, planningData)
            .pipe(catchError(this.handleError));
    }

    // =========================================================================
    // EXAMPLE MODEL METHODS
    // =========================================================================

    /**
     * Get list of available example models from backend.
     * Returns filenames like ["example1.pnml", "example2.pnml"]
     */
    getAvailableExampleModels(): Observable<string[]> {
        const url = `${environment.backendApiUrl}/simulation-petri-nets/example-models`;
        return this.http.get<string[]>(url)
            .pipe(catchError(this.handleError));
    }

    /**
     * Load a specific example model by name.
     * Returns the raw PNML content as a string.
     */
    loadExampleModelByName(name: string): Observable<string> {
        const url = `${environment.backendApiUrl}/simulation-petri-nets/example-models/${encodeURIComponent(name)}`;
        return this.http.get(url, { responseType: 'text' })
            .pipe(catchError(this.handleError));
    }

    // =========================================================================
    // ERROR HANDLING
    // =========================================================================

    /**
     * Centralized error handler for all HTTP requests.
     * Logs the error and returns a user-friendly Observable error.
     */
    private handleError(error: HttpErrorResponse): Observable<never> {
        let message: string;

        if (error.error instanceof ErrorEvent) {
            // Client-side error (network issues, etc.)
            message = `Network error: ${error.error.message}`;
        } else {
            // Server-side error
            message = `Server error: ${error.status} - ${error.message}`;
        }

        console.error('[PlanningService]', message);
        return throwError(() => new Error(message));
    }
}

// =========================================================================
// TYPE DEFINITIONS
// =========================================================================

interface SimulationResult {
    /** Sequence of transition IDs that fired */
    firing_seq: string[];
    /** Token counts at each step (optional) */
    detailed_log?: StepLog[];
    /** Statistics for multi-run simulations */
    statistics?: TransitionStats;
}

interface StepLog {
    step: number;
    transition: string;
    marking: { [placeId: string]: number };
}

interface TransitionStats {
    [transitionId: string]: {
        fire_count: number;
        fire_percentage: number;
    };
}

interface ParameterDefaults {
    [paramName: string]: number | string | boolean;
}

interface PlanningRequest {
    model_pnml: string;
    parameters: { [key: string]: any };
}

interface PlanningResult {
    schedule: any[];
    metrics: { [key: string]: number };
}
\end{lstlisting}

\section{Usage in Components}
\label{sec:planning-usage}

\subsection{Running a Simulation}

\begin{lstlisting}[language=TypeScript, caption={Using PlanningService in ButtonBarComponent}]
import { Component } from '@angular/core';
import { PlanningService } from '../../tr-services/planning.service';
import { PnmlService } from '../../tr-services/pnml.service';
import { UiService } from '../../tr-services/ui.service';

@Component({ /* ... */ })
export class ButtonBarComponent {
    constructor(
        private planningService: PlanningService,
        private pnmlService: PnmlService,
        private uiService: UiService
    ) {}

    /**
     * Called when user clicks "Run Simulation" button.
     */
    runSimulation(): void {
        // 1. Get current Petri net as PNML string
        const pnmlContent = this.pnmlService.getPNML();

        // 2. Send to backend for simulation
        this.planningService
            .runSimpleSimulationFromString(pnmlContent, 1)
            .subscribe({
                next: (result) => {
                    console.log('Simulation complete:', result.firing_seq.length, 'steps');
                    
                    // 3. Update UI with results
                    this.uiService.setSimulationResults(result);
                    this.uiService.tab = TabState.Simulation;
                },
                error: (err) => {
                    console.error('Simulation failed:', err.message);
                    this.showErrorDialog('Simulation failed: ' + err.message);
                }
            });
    }
}
\end{lstlisting}

\subsection{Loading Example Models}

\begin{lstlisting}[language=TypeScript, caption={Loading example models with dropdown}]
import { Component, OnInit } from '@angular/core';
import { Observable } from 'rxjs';
import { PlanningService } from '../../tr-services/planning.service';
import { PnmlService } from '../../tr-services/pnml.service';

@Component({
    selector: 'app-example-selector',
    template: `
        <select (change)="onModelSelected($event)">
            <option value="">-- Select Example --</option>
            <option *ngFor="let model of availableModels$ | async" [value]="model">
                {{ model }}
            </option>
        </select>
    `
})
export class ExampleSelectorComponent implements OnInit {
    availableModels$: Observable<string[]>;

    constructor(
        private planningService: PlanningService,
        private pnmlService: PnmlService
    ) {}

    ngOnInit(): void {
        // Fetch available examples on component init
        this.availableModels$ = this.planningService.getAvailableExampleModels();
    }

    onModelSelected(event: Event): void {
        const select = event.target as HTMLSelectElement;
        const modelName = select.value;

        if (!modelName) return;

        // Load the selected model
        this.planningService.loadExampleModelByName(modelName).subscribe({
            next: (pnmlContent) => {
                // Parse and load into DataService
                this.pnmlService.parse(pnmlContent);
                console.log('Loaded example:', modelName);
            },
            error: (err) => {
                console.error('Failed to load example:', err.message);
            }
        });
    }
}
\end{lstlisting}

\section{Request/Response Flow}
\label{sec:planning-flow}

\begin{figure}[htb]
    \centering
    \begin{tikzpicture}[
        node distance=1.2cm,
        box/.style={rectangle, draw, minimum width=2.2cm, minimum height=0.6cm, font=\small},
    ]
        % Actors
        \node[box] (user) at (0,0) {User};
        \node[box] (btn) at (3,0) {ButtonBar};
        \node[box] (pnml) at (6,0) {PnmlService};
        \node[box] (plan) at (9,0) {PlanningService};
        \node[box] (back) at (12,0) {Backend};
        
        % Flow
        \draw[->] (user) -- ++(0,-1) -| node[pos=0.25, below] {\tiny clicks Run} (btn);
        \draw[->] (btn) -- ++(0,-2) -| node[pos=0.25, below] {\tiny getPNML()} (pnml);
        \draw[->] (pnml) -- ++(0,-3) -| node[pos=0.25, below] {\tiny PNML string} (btn);
        \draw[->] (btn) -- ++(0,-4) -| node[pos=0.25, below] {\tiny runSimulation()} (plan);
        \draw[->] (plan) -- ++(0,-5) -| node[pos=0.25, below] {\tiny POST /simple-sim} (back);
        \draw[->] (back) -- ++(0,-6) -| node[pos=0.25, below] {\tiny JSON result} (plan);
        \draw[->] (plan) -- ++(0,-7) -| node[pos=0.25, below] {\tiny Observable<Result>} (btn);
        
    \end{tikzpicture}
    \caption{Simulation request flow through PlanningService}
    \label{fig:simulation-flow}
\end{figure}

\section{Backend API Endpoints}
\label{sec:planning-endpoints}

The backend exposes the following REST endpoints:

\begin{table}[htb]
    \centering
    \begin{tabular}{llp{6cm}}
        \toprule
        \textbf{Method} & \textbf{Endpoint} & \textbf{Description} \\
        \midrule
        GET & \texttt{/dto/planning/defaults} & Get default parameter values \\
        POST & \texttt{/dto/planning} & Run planning optimization \\
        POST & \texttt{/simulation-petri-nets/simple-sim} & Run Petri net simulation \\
        GET & \texttt{/simulation-petri-nets/example-models} & List available examples \\
        GET & \texttt{/simulation-petri-nets/example-models/:name} & Get example PNML content \\
        \bottomrule
    \end{tabular}
    \caption{Backend API endpoints}
    \label{tab:api-endpoints}
\end{table}

\section{Error Handling Patterns}
\label{sec:planning-errors}

Robust error handling ensures graceful degradation when the backend is unavailable:

\begin{lstlisting}[language=TypeScript, caption={Error handling with retry and fallback}]
import { catchError, retry, timeout } from 'rxjs/operators';
import { of } from 'rxjs';

// With retry and timeout
this.planningService.runSimpleSimulation(file)
    .pipe(
        timeout(30000),  // 30 second timeout
        retry(2),        // Retry twice on failure
        catchError(err => {
            // Fallback to empty result
            console.warn('Simulation failed, using empty result');
            return of({ firing_seq: [], detailed_log: [] });
        })
    )
    .subscribe(result => {
        // Handle result (real or fallback)
    });
\end{lstlisting}

\section{CORS Configuration}
\label{sec:planning-cors}

For development, the backend must allow cross-origin requests from the Angular dev server:

\begin{lstlisting}[language=Python, caption={Flask CORS configuration (backend)}]
from flask import Flask
from flask_cors import CORS

app = Flask(__name__)
CORS(app, resources={
    r"/l3s-offshore-2/*": {
        "origins": ["http://localhost:4200"],
        "methods": ["GET", "POST", "OPTIONS"],
        "allow_headers": ["Content-Type"]
    }
})
\end{lstlisting}

In production, the nginx reverse proxy handles CORS headers.

\section{Summary}
\label{sec:planning-summary}

The \texttt{PlanningService} extends frontend capabilities through backend communication:

\begin{itemize}
    \item \textbf{HTTP Communication}: Uses Angular's \texttt{HttpClient} for REST API calls
    \item \textbf{Environment Configuration}: Separate URLs for development and production
    \item \textbf{File Uploads}: Uses \texttt{FormData} to send PNML files
    \item \textbf{Error Handling}: Centralized handler with user-friendly messages
    \item \textbf{Observable Pattern}: Async results via RxJS Observables
\end{itemize}

This service enables the application to leverage powerful server-side algorithms while maintaining a responsive, client-side user interface.
